\subsection{EXAMPLES}

Example 1. Let \(\mathfrak{g}\) be a Lie algebra over \(\mathbf{K}\) and \(\mathbf{M}\) a \(\mathfrak{g}\) -module. The \(\mathfrak{g}\) -module structure on \(\mathbf{M}\) and the trivial \(\mathfrak{g}\) -module structure on \(\mathbf{K}\) define a \(\mathfrak{g}\) -module structure on the \(\mathbf{K}\) -module \(N = \mathcal{L}(M, M; K)\) of bilinear forms on \(\mathbf{M}\) . Then

\[
(x _ {N}. \beta) (m, m ^ {\prime}) = - \beta (x _ {M}. m, m ^ {\prime}) + \beta (m, x _ {M}. m ^ {\prime})\tag{10}
\]

for all \(x \in \mathfrak{g}\) , \(m, m'\) in \(\mathbf{M}\) , \(\beta \in \mathbf{N}\) . If \(\beta\) is a given element of \(\mathbf{N}\) , the set of \(x \in \mathfrak{g}\) such that \(x_{\mathbf{N}}. \beta = 0\) is a subalgebra of \(\mathfrak{g}\) .

Let \(\mathbf{M}\) be a K-module and \(\beta\) a bilinear form on \(\mathbf{M}\) . By the above, the set of \(x \in \mathfrak{gl}(\mathbf{M})\) such that

\[
\beta (x. m, m ^ {\prime}) + \beta (m, x. m ^ {\prime}) = 0
\]

for all \(m \in \mathbf{M}\) and \(m' \in \mathbf{M}\) is a Lie subalgebra of \(\mathfrak{gl}(\mathbf{M})\) . Suppose that \(\mathbf{K}\) is a field, \(\mathbf{M}\) is finite-dimensional and \(\beta\) is non-degenerate. Then every \(x \in \mathfrak{gl}(\mathbf{M})\) admits a left adjoint \(x^*\) (relative to \(\beta\) ) which is everywhere defined and the subalgebra in question is the set of \(x \in \mathfrak{gl}(\mathbf{M})\) such that \(x^* = -x\) . By this process we can construct two important examples of Lie algebras:

\begin{enumerate}
\def\labelenumi{(\alph{enumi})}
\tightlist
\item
  Take \(\mathbf{M} = \mathbf{K}^n\) and
\end{enumerate}

\[
\beta ((\xi_ {1}, \dots , \xi_ {n}), (\eta_ {1}, \dots , \eta_ {n})) = \xi_ {1} \eta_ {1} + \dots + \xi_ {n} \eta_ {n}.
\]

We canonically identify \(\mathfrak{gl}(\mathbf{K}^n)\) with \(\mathbf{M}_n(\mathbf{K})\) . Then the Lie algebra obtained is the Lie algebra of skew-symmetric matrices. \(*\) (When \(\mathbf{K} = \mathbf{R}\) , this algebra is the Lie algebra of the orthogonal group \(\mathbf{O}(n, \mathbf{R}))\) .

\begin{enumerate}
\def\labelenumi{(\alph{enumi})}
\setcounter{enumi}{1}
\tightlist
\item
  Take \(\mathbf{M} = \mathbf{K}^{2m}\) and
\end{enumerate}

\[
\beta ((\xi_ {1}, \dots , \xi_ {2 m}), (\eta_ {1}, \dots , \eta_ {2 m})) = \xi_ {1} \eta_ {m + 1} - \eta_ {1} \xi_ {m + 1} + \dots + \xi_ {m} \eta_ {2 m} - \eta_ {m} \xi_ {2 m}.
\]

The matrix of \(\beta\) with respect to the canonical basis of \(\mathbf{K}^{2m}\) is the matrix \(\left( \begin{array}{cc}0 & I_m\\ -I_m & 0 \end{array} \right)\) . Let \(U = \left( \begin{array}{cc}A & B\\ C & D \end{array} \right)\) be the matrix with respect to the canonical basis of \(\mathrm{K}^{2m}\) of an element \(u\) of \(\mathfrak{gl}(\mathbf{M})\) ( \(A,B,C,D\) lying in \(\mathbf{M}_m(\mathbf{K})\) ). By formula (50) of Algebra, Chapter IX, § 1, no. 10, \(u^*\) has with respect to the same basis the matrix

\[
\left( \begin{array}{c c} 0 & - I _ {m} \\ I _ {m} & 0 \end{array} \right) \left( \begin{array}{c c} ^ {t} A & ^ {t} C \\ ^ {t} B & ^ {t} D \end{array} \right) \left( \begin{array}{c c} 0 & I _ {m} \\ - I _ {m} & 0 \end{array} \right) = \left( \begin{array}{c c} ^ {t} D & - ^ {t} B \\ - ^ {t} C & ^ {t} A \end{array} \right).
\]

The condition \(u^{*} = -u\) is therefore equivalent to the conditions

\[
D = - ^ {t} A \qquad B = ^ {t} B \qquad C = ^ {t} C.
\]

*When \(\mathbf{K} = \mathbf{R}\) , the Lie algebra obtained is the Lie algebra of the symplectic group \(\mathbf{Sp}(2m, \mathbf{R})\) .*

Example 2. We preserve the notation of Example 1.

The g-module structure on M defines on the K-module \(P = \mathcal{L}_{K}(M, M)\) of endomorphisms of M a g-module structure. By (6), for all \(x \in g\) and \(u \in P\) :

\[
x _ {\mathrm{P}}. u = \left[ x _ {\mathrm{M}}, u \right] = (\operatorname{ad} x _ {\mathrm{M}}). u\tag{11}
\]

where ad \(x_{M}\) denotes the image of \(x_{M}\) under the adjoint representation of \(\mathfrak{gl}(M)\) . In other words:

\[
x _ {\mathrm{P}} = \operatorname{ad} x _ {\mathrm{M}}\tag{12}
\]

in \(\mathcal{L}(\mathcal{L}(\mathrm{M},\mathrm{M})) = \mathcal{L}(\mathfrak{gl}(\mathrm{M}))\) .

